\begin{helper}[Density of a ray]
Let $F$ be a nontrivial front on $X$ and let $n\in X$.  Every infinite
$Y\subseteq\mathsf{TailSet}(X,n)$ has a member $s$ of $\mathsf{Ray}(F,n)$
with $\mathsf{ProperPrefixSet}(s,Y)$.
\end{helper}
\begin{proof}
By \noderef{FrontNontrivialBase}, every member of $F$ is nonempty.  Let
$Y\subseteq\mathsf{TailSet}(X,n)$ be infinite.  Then
$\{n\}\cup Y$ is infinite because it contains the infinite set $Y$, and
it is contained in $X$ because $n\in X$ and $Y\subseteq X$ by
\noderef{TailSet}.  Apply the density clause of \noderef{Front} to obtain
$u\in F$ and a cutoff $m\in\{n\}\cup Y$ such that

\[
u=\{k\in\{n\}\cup Y\mid k<m\}.
\]

The cutoff cannot be $n$.  Indeed, in that case the displayed set would
be empty: $n$ is not less than itself, and every element of $Y$ is greater
than $n$ by \noderef{TailSet}.  This contradicts the nonemptiness of $u$.
Thus $m\in Y$, and $n<m$.  The displayed equation now shows $n\in u$.

Set $s=u\setminus\{n\}$.  Every element of $s$ lies in $Y$ and is greater
than $n$, while inserting $n$ back gives $\{n\}\cup s=u\in F$.  Hence
$s\in\mathsf{Ray}(F,n)$ by \noderef{Ray}.  Removing $n$ from the displayed
cutoff equation gives

\[
s=\{k\in Y\mid k<m\},
\]

so $\mathsf{ProperPrefixSet}(s,Y)$ by \noderef{ProperPrefixSet}.  This is
the required ray-density conclusion.
\end{proof}
